\section{Introduction}
\label{sec:introduction}

Gene expression is regulated in a cell-type- and context-specific manner. The precise control of gene transcription is achieved by dynamic changes of chromatin accessibility at \textit{cis}-regulatory regions, underlying diverse cellular functions. Therefore, mapping enhancers, an important class of gene regulatory elements, to their target genes can provide insights into the normal physiology and complex human diseases \cite{gasperini2020enhancers}. 
Advances in single-cell RNA sequencing (scRNA-seq) and the single-cell assay for transposase-accessible chromatin using sequencing (scATAC-seq) have enabled the profiling of gene expression and chromatin accessibility at single cell level. Methods based on the high-throughput sequencing enable the mapping of expression quantitative trait loci (eQTLs) and chromatin accessibility quantitative trait loci (caQTLs) at cell type resolution. These advances have motivated the development of statistical methods for linking candidate enhancers to their target genes in specific cellular contexts. For example, SHARE-seq \cite{ma2020shareseq} and Signac \cite{stuart2021signac} infer putative \textit{cis}-regulatory links by correlating chromatin accessibility with gene expression across cells. scMultiMap \cite{su2025scmultimap} jointly models RNA and ATAC counts and estimates enhancer--gene correlations. pgBoost \cite{dorans2025pgboost} integrates linking scores from multiple single-cell multiome methods with genomic distance to prioritize candidate variant--gene links. However, most existing approaches rely on marginal associations across cells or individuals, leaving the causal effect of peak accessibility on gene expression unresolved.

\noindent Context-specific enhancer--gene regulation may be biased or masked in marginal associations. Theoretically, the marginal association between peak accessibility and gene expression equals the causal effect of accessibility on expression plus a bias term contributed by any unmeasured factor that affects both, such as cell identity, transcription factor (TF) activity, or donor-level variation (\nameref{S1_Text}, Section~1). Depending on the confounding effects, a marginal association can arise in the absence of any causal effect, whereas a true causal effect can be masked, biased, or even reversed in direction. For example, in an analysis of 38 mouse cortical neuronal subtypes, \textit{Stim2} expression was strongly correlated with the accessibility of a nearby enhancer ($r = 0.87$) but also with 23,526 enhancers located on other chromosomes, most of which likely reflect shared upstream regulation, a confounding factor that drives both the accessibility of these enhancers and the expression of \textit{Stim2} across subtypes, rather than direct regulatory interactions \cite{xie2023robust}. In addition, in the TRIPOD study, a TF that was required for expression of the target gene was negatively correlated with peak accessibility across cells; acting as a confounding factor, it offset the effect of the peak, so that the marginal peak--gene association was no longer significant \cite{jiang2022tripod}. Similarly, coordinated variation across donors can induce spurious associations: if mean accessibility and mean expression are both high in one donor and both low in another, a peak--gene pair that is unassociated within each donor may appear associated once cells from the two donors are pooled \cite{su2025scmultimap}. These examples suggest that marginal association does not necessarily imply causation. In this paper, we focus on the identification of causal enhancer--gene regulation.

\noindent Based on Mendel's laws of inheritance, Mendelian randomization (MR) uses genetic variants, usually single-nucleotide polymorphisms (SNPs), as instrumental variables to estimate the causal effect of exposures on outcomes even in the presence of unmeasured confounding \cite{daveysmith2003mendelian,lawlor2008mendelian}. The availability of previously published summary statistics from large-scale genetic association studies has made MR a practical and widely applicable approach for investigating the causal effects of exposures, such as body mass index (BMI) or gene expression, on health outcomes, such as complex diseases \cite{burgess2013summarized,holmes2014bmi,zhu2016smr,richardson2020twmr}. However, few studies have applied MR to transcriptional regulation where both the exposure and the outcome are molecular traits. In this paper, we use summary statistics from \textit{cis}-caQTL and \textit{cis}-eQTL studies and select genetic variants associated with chromatin accessibility as instrumental variables to estimate the causal effect of peak accessibility on gene expression.

\noindent Extending MR to study enhancer--gene regulation using \textit{cis}-QTL summary statistics presents unique challenges. First, standard MR methods were developed for complex-trait exposures, which typically have many independent genome-wide significant instruments, whereas \textit{cis}-caQTLs for a given peak are limited in number and confined to a local genomic region. Because these \textit{cis}-caQTLs are frequently in linkage disequilibrium (LD) with one another \cite{kang2026fusiomr,lin2024cismrcml}, stringent LD clumping to enforce independence can leave too few instrumental variables for analysis. Second, horizontal pleiotropy is pervasive in genetics: a genetic variant may affect target gene expression through other regulatory elements, violating the exclusion-restriction assumption \cite{bowden2015egger,kang2026fusiomr}. Third, because different MR estimators rely on distinct modeling assumptions, causal conclusions drawn from a single MR method may not be robust \cite{burgess2017interpreting}.

\noindent To address these challenges, we present a flexible MR framework to infer causal enhancer--gene regulation using population-scale caQTL and eQTL summary statistics. The framework comprises four methods: causal inverse-variance weighting (causal IVW), causal MR-Egger (causal Egger), causal summary-data-based Mendelian randomization (causal SMR), and causal joint. Both causal IVW and causal Egger use generalized least squares to account for correlations among genetic instruments \cite{burgess2016combining,burgess2017interpreting}, with causal IVW allowing for balanced pleiotropy and causal Egger allowing for directional pleiotropy \cite{burgess2017interpreting}. In contrast, causal SMR uses the lead caQTL for each peak as its sole instrument \cite{zhu2016smr}. Finally, causal joint combines the $p$-values from these three methods using the Cauchy combination test \cite{liu2020cauchy}, aiming to improve both robustness and statistical power. To evaluate our MR framework, we benchmark it against two marginal methods for identifying enhancer--gene links, computing peak--gene correlations across metacells (marginal multiomics) and across individuals within each cell type (marginal pseudobulk).

\noindent We applied our MR framework to cell-type-specific \textit{cis}-caQTL and \textit{cis}-eQTL summary statistics from the Chinese Immune Multi-Omics Atlas (CIMA) \cite{yin2026cima}, focusing on three lymphocyte populations: naive B cells (Bn cells), naive CD4$^+$ T cells (CD4 T cells), and naive CD8$^+$ T cells (CD8 T cells). We found that the MR estimates differed substantially from marginal associations. Although the two approaches showed concordance on many peak--gene pairs, they differed in effect size or direction for others. Notably, our MR framework uncovered enhancer--gene links missed by both marginal methods, referred to as MR-only links. These MR-only links were well supported by orthogonal regulatory resources, including enhancer databases, \textit{cis}-eQTL studies, promoter capture Hi-C, and H3K27ac HiChIP data. Moreover, compared with marginal methods, MR-inferred links exhibited higher odds ratios for overlap with the enhancer-driven regulons (eRegulons) reported in CIMA. To illustrate the utility of our framework for interpreting genome-wide association study (GWAS) variants in disease-relevant cell types, we evaluated heritability enrichment across seven immune-mediated diseases. Peaks involved in MR-only links showed higher heritability enrichment than all candidate peaks, and our framework further pinpointed enhancer--gene links harboring disease-associated GWAS variants, offering insights into their regulatory mechanisms. Together, these analyses demonstrate how population-scale single-cell multiomics and causal inference can be combined to dissect cell-type-specific enhancer--gene regulation.
